% !TEX root = ../../main.tex
%% Défi du chapitre 11 — La fonction W de Lambert
%% Résultats vérifiés par verifications/defis_11_15.py.

\rubrique{Défi}

\begin{defi}{La fonction W de Lambert}{2 h}
L'équation $x\,\mathrm{e}^{x}=y$ n'a pas de solution « exprimable » à l'aide
des fonctions usuelles ; on en fait donc une fonction nouvelle, la
fonction $W$ de Lambert, omniprésente en physique, en combinatoire et en
informatique. On l'étudie ici entièrement par l'analyse asymptotique.

\partieExo{A}{Définition et dérivée}
\begin{questions}
  \item Étudier $g:x\mapsto x\,\mathrm{e}^{x}$ sur $\intFO{-1}{+\infty}$ et
    montrer qu'elle réalise une bijection de $\intFO{-1}{+\infty}$ sur
    $\intFO{-\frac1{\mathrm{e}}}{+\infty}$. On note $W$ sa réciproque.
    Préciser $W(0)$, $W(\mathrm{e})$ et $W\bigl(-\frac1{\mathrm{e}}\bigr)$.
  \item Montrer que $W$ est continue et strictement croissante, dérivable
    sur $\intOO{-\frac1{\mathrm{e}}}{+\infty}$, avec
    \[
    W'(y)=\frac{1}{\bigl(1+W(y)\bigr)\mathrm{e}^{W(y)}}
    \qquad\text{et}\qquad
    W'(y)=\frac{W(y)}{y\bigl(1+W(y)\bigr)}\ \ (y\neq0).
    \]
    Que vaut $W'(0)$ ? Que se passe-t-il en $-\frac1{\mathrm{e}}$ ?
\end{questions}

\medskip
\begin{center}
\begin{tikzpicture}[scale=1.35,line cap=round]
  \draw[otRule,line width=0.3pt] (-1.2,-1.2) grid (2.3,2.3);
  \draw[otMuted,line width=0.5pt,-{Stealth[length=1.6mm]}] (-1.3,0) -- (2.45,0);
  \draw[otMuted,line width=0.5pt,-{Stealth[length=1.6mm]}] (0,-1.3) -- (0,2.45);
  \draw[otMuted,line width=0.5pt,dash pattern=on 2pt off 1.6pt] (-1.2,-1.2) -- (2.3,2.3);
  \draw[otNavy,line width=1.2pt] plot[domain=-1:0.9,samples=100] ({\x},{\x*exp(\x)});
  \draw[otGreenDark,line width=1.2pt] plot[domain=-1:0.9,samples=100] ({\x*exp(\x)},{\x});
  \fill[otNavy] (-1,-0.3679) circle (1.4pt);
  \fill[otGreenDark] (-0.3679,-1) circle (1.4pt);
  \node[font=\scriptsize,text=otNavy,anchor=south west] at (0.75,1.9) {$y=x\,\mathrm{e}^{x}$};
  \node[font=\scriptsize,text=otGreenDark,anchor=west] at (2.05,0.75) {$W$};
  \node[font=\scriptsize,text=otMuted,anchor=north west] at (-0.36,-1) {$\bigl(-\tfrac1{\mathrm{e}},-1\bigr)$};
  \node[font=\tiny,text=otMuted,below] at (1,0) {$1$};
  \node[font=\tiny,text=otMuted,left] at (0,1) {$1$};
\end{tikzpicture}\\[0.5mm]
\legende{La fonction $x\mapsto x\,\mathrm{e}^{x}$ (en bleu) et sa réciproque
$W$ (en vert), symétriques par rapport à la première bissectrice ; $W$ a une
tangente verticale au point $\bigl(-\frac1{\mathrm{e}},-1\bigr)$.}
\end{center}

\partieExo{B}{Développement limité en $0$}
On admet que $W$ admet en $0$ un développement limité à tout ordre.
\begin{questions}[resume]
  \item En identifiant les coefficients dans $W(y)\,\mathrm{e}^{W(y)}=y$,
    montrer que $W(y)=y-y^{2}+\dfrac32y^{3}+o\bigl(y^{3}\bigr)$.
  \item Pousser le calcul à l'ordre $4$ ; on vérifiera que le coefficient
    de $y^{n}$ vaut $(-n)^{n-1}/n!$ pour $n\leqslant4$ (c'est vrai pour
    tout $n$ : formule d'inversion de Lagrange, admise).
  \item En déduire une valeur approchée de la solution de
    $x\,\mathrm{e}^{x}=0{,}01$, en précisant l'ordre de grandeur de
    l'erreur.
\end{questions}

\partieExo{C}{Comportement en $+\infty$}
\begin{questions}[resume]
  \item Montrer que $W(y)\to+\infty$ quand $y\to+\infty$, que
    $W(y)+\ln W(y)=\ln y$ pour $y>0$, et en déduire $W(y)\sim\ln y$.
  \item Établir successivement
    \[
    W(y)=\ln y-\ln\ln y+o(1),
    \qquad
    W(y)=\ln y-\ln\ln y+\frac{\ln\ln y}{\ln y}
    +o\Bigl(\frac{\ln\ln y}{\ln y}\Bigr).
    \]
  \item Pour $y>1$, montrer que l'équation $x^{x}=y$ admet une unique
    solution $x_y>1$, que
    \[
    x_y=\frac{\ln y}{W(\ln y)},
    \]
    et en déduire un équivalent simple de $x_y$ quand $y\to+\infty$.
\end{questions}
\end{defi}
